\pagenumbering{roman}

% ---------- title page -------------------------------------------------------
\begin{titlepage}
\centering
\vspace*{1.2cm}
{\Large University College London\par}
\vspace{0.3cm}
{\large Department of Computer Science\par}
\vfill
{\LARGE\bfseries When Is Expensive Perception Worth Paying For?\par}
\vspace{0.5cm}
{\large Measuring the Ceiling, the Predictability, and the Economics\\
of Representation Routing in Web Agents\par}
\vfill
{\large Jiaming Wei\par}
\vspace{0.4cm}
{\normalsize MSc Artificial Intelligence for Sustainable Development\par}
{\normalsize COMP0191: MSc Project\par}
\vfill
{\normalsize September 2026\par}
\vspace{0.8cm}
{\footnotesize
This report is submitted as part requirement for the MSc degree in Artificial
Intelligence for Sustainable Development at University College London. It is
substantially the result of my own work except where explicitly indicated in
the text. The report may be freely copied and distributed provided the source
is explicitly acknowledged.\par}
\end{titlepage}

% ---------- abstract ---------------------------------------------------------
\chapter*{Abstract}
\addcontentsline{toc}{chapter}{Abstract}
\markboth{Abstract}{Abstract}

A web agent looks at a web page, decides what to click or type, acts, and
repeats until a user's task is done. Before it can act it must be handed some
encoding of the page: cheap text listing the page's elements and their labels,
an expensive screenshot with the interactable elements numbered and boxed on
it, or both together. The choice is normally made once, when the system is configured,
and applied everywhere. Since a task runs for tens of steps, it is also paid for
tens of times. This dissertation asks whether the expensive choice is necessary
at every step, and whether the steps where it \emph{is} necessary can be
identified cheaply enough for the identification to be worth making.

We hold the agent fixed and vary only what it is shown. Six \emph{observation
modes} differ in the page text supplied, the wording of the instructions around
it, and whether an annotated image is sent at all; the task, the available
actions, the step limit and the rules for what counts as a cost are identical
across all six. We run all six on 8 website-and-model combinations, drawn from
two public benchmarks of website tasks --- VisualWebArena and WebArena --- and
three underlying models. On that base we measure three quantities in sequence:
what an ideal after-the-fact selector --- an \emph{oracle} allowed to pick, for
each task, a mode already known to have solved it --- would gain by choosing per
task; whether the same choice can be made in advance from information a deployed
system would have; and whether acting on those predictions beats always using
the cheapest mode, once cost is counted.

The answers are, in order, \emph{yes}, \emph{no}, and \emph{the question is
mis-posed}. The oracle gains $+3.45$ to $+16.35$ percentage points of task
success over the best single mode while spending 1.6--35.3\% less, in every
combination and on both benchmarks. Our own noise standard deflates that
headline: running one unchanged mode a second time buys 2.0--7.6 points by
itself, and when both sides are allowed six runs, six repetitions of the
strongest single mode come within 3.53 points of the six-mode oracle --- below
the threshold we derive for what a single repetition can produce on its own.
What survives is the ceiling's price rather than its height, since selection
reaches that level with one attempt per task where repetition needs six. A
policy wins here only if it is \emph{both} more successful and cheaper than the
one it is measured against --- a \emph{Pareto} improvement. Under a fully nested cross-validation protocol ---
every modelling choice made on training tasks alone, then tested on tasks held
back --- the learned router --- which only has to predict whether \emph{any}
mode will solve the task, not which of the six to use --- clears that bar
against always-cheapest in \textbf{0 of 8} combinations. The oracle, told the answers in
advance, fails to clear it in 7 of 8. That recasts the negative result from
``the learner failed'' to ``no Pareto-improving operating point exists against
always-cheapest under this cost boundary''.

What blocks the prediction is the supply of training examples, not the capacity
of the model. A task can teach the router something only if some mode solved it,
so at 2--27\% base success 4 of 6 combinations lack enough solved tasks per mode
to train and test a classifier at all. The input carrying most of the apparent
predictive power on VisualWebArena turns out to be the benchmark's own
human-written difficulty annotation, which no deployment has. Training on
deliberately reduced data confirms that scarcity is part of the story and then
bounds it: the learning curves are still rising, so more data would buy a better
predictor, but the best any predictor can become is the oracle, and the oracle
wins in 1 of 8 combinations. Pricing this constraint puts the benchmark that
would relieve it at 2.1 to 4.2 times the size of the ones that exist, which is a
specification rather than an impossibility.

The contribution is not a router. It is a measured boundary: the conditions
under which \emph{representation routing} --- choosing the page encoding per
task --- cannot be learned, with the
opportunity proven to exist before it is declared unreachable, and the failure
traced to a mechanism we can price.

% ---------- acknowledgements -------------------------------------------------
\chapter*{Acknowledgements}
\addcontentsline{toc}{chapter}{Acknowledgements}
\markboth{Acknowledgements}{Acknowledgements}

I thank my supervisor for consistently redirecting attention from what the
system could be made to do towards what the evidence could actually support,
and Zekun Wu for the structural writing guidance that shaped how this document
separates what was measured from what was inferred. Compute for the local
backbones was provided by the Holistic AI GB10 nodes and by a dedicated A100
instance; neither was shared during the paper-grade runs, which matters for the
latency and energy figures reported in Chapter~\ref{ch:system}.

% ---------- public access ----------------------------------------------------
\chapter*{Data, Code and Public Access}
\addcontentsline{toc}{chapter}{Data, Code and Public Access}
\markboth{Data, Code and Public Access}{Data, Code and Public Access}

Every component of this study is publicly obtainable with one stated exception.
The benchmarks (VisualWebArena, WebArena) are open-source and self-hosted from
their published Docker images; two of the three backbones
(\texttt{Qwen3-VL-4B}, \texttt{google/gemma-3-4b-it}) are openly available model
weights run locally; the experiment harness, analysis pipeline and figure
scripts are released with the submission.

\textbf{The exception is the large backbone.} \texttt{Qwen3-VL-235B-A22B} was
served through a metered API proxy rather than self-hosted. Its outputs,
per-step logs and billing records are archived and released, but a third party
cannot re-run those conditions without equivalent proxy access. Results that
depend on that backbone alone are therefore \emph{replicable from our logs} but
not \emph{independently reproducible end to end}. Every claim that rests on
that backbone is marked where it is made, and
Chapters~\ref{ch:repr}--\ref{ch:e2e} are built so that none of the conclusions
depends on a single backbone.

\clearpage
\tableofcontents
\listoffigures
\listoftables
\clearpage

% ---------- reader's guide ---------------------------------------------------
\chapter*{Reader's Guide: Modes, Codes and Terms}
\addcontentsline{toc}{chapter}{Reader's Guide: Modes, Codes and Terms}
\markboth{Reader's Guide}{Reader's Guide}

This dissertation compares six ways of showing a web page to the same program,
across eight combinations of website and model. Those six ways and eight
combinations are referred to by short codes throughout. This page decodes them,
and defines the recurring terms in plain language. Nothing here is needed to
follow the argument of Chapter~\ref{ch:intro}; it is placed first so that the
tables and figures in later chapters can be read without hunting for a
definition.

\paragraph{The six observation modes.} Each mode is one complete recipe for what
the program is shown at a single browser step. They differ in three things: the
page text supplied, the wording of the instructions around that text, and
whether an image is sent.

\begin{center}\small
\begin{tabularx}{\textwidth}{@{}lX@{}}
\toprule
Code & What the program is shown \\
\midrule
\mDOM{}  & A structured text listing of the page's elements, their labels and
           their roles. No image. \\
\mPP{}   & The same text listing, but with instructions written as though a
           numbered screenshot were coming. No image. \\
\mPT{}   & Only the lines of that listing that carry an element number, with
           the structure flattened away. No image. \\
\mPS{}   & Those same numbered lines, with instructions written as though a
           numbered screenshot were coming. No image. \\
\mSoM{}  & Those same numbered lines \emph{and} a screenshot with the
interactable elements numbered and boxed on it. \\
\mVis{}  & The plain screenshot only. No page text at all. \\
\bottomrule
\end{tabularx}
\end{center}

\noindent The four modes with no image are the cheap end; \mSoM{} is the
expensive end. \mPT{}, \mPP{} and \mPS{} exist to separate three changes that
happen together when one moves from \mDOM{} to \mSoM{}, and are called
\emph{phantom} modes throughout (Section~\ref{sec:modes}).

\paragraph{The eight website-and-model combinations.} A \emph{cell} is one
website paired with one model. Every cell runs all six modes over the same
tasks.

\begin{center}\small
\begin{tabularx}{\textwidth}{@{}llX@{}}
\toprule
Code & Stands for & Which is \\
\midrule
\cell{cls}    & VisualWebArena classifieds & a self-hosted classified-ads site \\
\cell{red}    & VisualWebArena reddit      & a self-hosted discussion-forum site \\
\cell{wa\_red} & WebArena reddit           & the same forum in a second, text-oriented benchmark \\
\midrule
\cell{B0} & \texttt{Qwen3-VL-235B-A22B} & the large model, reached through a paid API \\
\cell{B1} & \texttt{Qwen3-VL-4B}        & a small model of the same design family, run locally \\
\cell{B2} & \texttt{google/gemma-3-4b-it} & a small model of a \emph{different} design family, run locally \\
\bottomrule
\end{tabularx}
\end{center}

\noindent So \cell{cls$\cdot$B0} means ``the classified-ads site, run with the
large model''. The eight cells are the three sites crossed with the models
available on each.

\paragraph{Recurring terms.}

\begin{center}\small
\begin{tabularx}{\textwidth}{@{}lX@{}}
\toprule
Term & Meaning in this dissertation \\
\midrule
step        & One look-decide-act cycle: the program is shown the page, picks one
              action, and the browser performs it. \\
episode     & One attempt at one task, from the starting page until it succeeds,
              gives up, errors, or runs out of steps. \\
task success & Whether the benchmark's own automatic checker judged the task
              complete. Success rates here are percentages of tasks. \\
backbone    & The underlying model doing the deciding (\cell{B0}, \cell{B1} or
              \cell{B2}). \\
token       & The unit in which model input is counted and billed; longer page
              text and larger images mean more tokens. \\
pp          & Percentage points. A move from 14\% to 17\% is $+3$ pp. \\
oracle & An idealised chooser that is told, after the fact, which modes solved
a task, and is allowed to pick one of them --- the cheapest, where more than
one worked. It cannot be built; it measures how much opportunity exists. \\
ceiling & How much better that oracle does than the best single mode. It is an
upper bound on what any chooser could \emph{gain}, not a result any system
attains. \\
drop-one contribution & How far the ceiling falls if one mode is deleted from
the oracle's menu. Running one mode a second time produces a number of the
same shape, so a positive value alone does not establish that the mode was
irreplaceable --- Chapter~\ref{ch:repr} downgrades an earlier claim on exactly
this point. \\
arm & One run of one mode over the whole task set. Six different modes run
once is six arms; one mode run six times is also six arms. Comparing a six-arm
result against a one-arm result would flatter whichever side had more
attempts, so this dissertation matches the count before comparing. \\ noise
floor & How much a result moves when an unchanged setup is simply run again.
Two versions are needed and they must not be mixed: one for claims about the
ceiling (what one more run buys), one for claims comparing two modes' success
rates. A single observed rerun difference is one draw rather than a threshold,
so the usable thresholds are larger than the raw differences: 4.91--7.59\pp{}
for ceiling claims and 3.82--4.15\pp{} for mode-to-mode comparisons.
Appendix~\ref{app:noise} derives both and states which comparisons each one
licenses. \\
routing     & Choosing which mode to use, rather than fixing one in advance.
              The router picks the mode; the program still picks the click. \\
triage & The easier half of routing: predicting whether \emph{any} mode will
solve this task at all. If none will, the cheapest is used, since nothing is
gained by paying more. Its positive class is therefore the set of tasks that
some mode solves --- which is also the only set that can say \emph{which} mode
to prefer. \\
fixed policy & Using the same mode for everything. \emph{Always-cheapest} is
the fixed policy that always uses the cheapest mode for that cell --- which
mode that is differs from cell to cell. \\
Pareto win  & Beating a comparison policy on \emph{both} counts at once: higher
              task success \emph{and} lower cost. Being better on one and worse
              on the other does not count. \\
serving-time information & What a deployed system would know before it has paid
              for anything: the task wording and the cheap page text. Excludes
              anything only knowable after the task is attempted, and anything
              the benchmark supplies that a deployment would not have, such as
              its own human-written difficulty rating. \\
nested cross-validation & A testing procedure in which the data is split
              repeatedly, and every choice --- which mode, which threshold, how
              to scale the inputs --- is made using only the training part, then
              judged on the part held back. It exists to stop a method from
              being tuned on the data it is then scored on. \\
AUROC & A single number between 0 and 1 summarising how well a predictor
\emph{ranks} one class above the other. 0.5 is chance; 1.0 is perfect. It says
nothing about whether acting on that ranking pays, once the two kinds of
mistake carry different costs --- Chapters~\ref{ch:pred} and~\ref{ch:e2e} turn
on exactly that gap. \\
\bottomrule
\end{tabularx}
\end{center}

\clearpage
\pagenumbering{arabic}
